% !TEX root = ../../main.tex
\subsection{威胁模型}

FoodShield 的威胁模型围绕外卖订单通信场景中的身份伪造、隐私泄露、日志篡改、管理员越权和匿名身份滥用等问题展开。系统假设攻击者具有一定的网络访问能力和业务接口调用能力，可能通过伪造请求参数、重放认证凭证、篡改数据库记录、批量注册账号或越权访问接口等方式破坏系统安全性。

本系统重点考虑以下几类威胁：

\begin{enumerate}[label=\arabic*., leftmargin=2em]
\item \textbf{订单身份伪造攻击}：攻击者可能尝试构造虚假的订单号或 Token，以非订单参与方身份进入用户与骑手之间的订单通信通道。如果系统仅依赖订单号作为入口，则订单号一旦泄露，攻击者可能冒充合法用户或骑手参与通信。

\item \textbf{Token 重放攻击}：攻击者可能截获某次合法订单认证请求中的 Token，并在之后重复使用该 Token 尝试进入通信通道。如果 Token 不包含时间戳、随机数或有效期限制，则存在被重复利用的风险。

\item \textbf{机器批量注册攻击}：攻击者可能使用脚本批量注册账号，生成大量匿名身份并发起垃圾订单或骚扰通信。如果注册阶段仅输入用户名并直接生成 PID，则系统实际上只有标识生成过程，而缺少有效认证和反滥用能力。

\item \textbf{用户身份暴露风险}：在普通外卖通信模式下，用户真实身份、联系方式、地址备注等信息可能被骑手端或管理员端过度接触。恶意骑手可能利用固定身份标识进行跨订单关联，甚至造成骚扰或隐私泄露。

\item \textbf{通信日志篡改攻击}：攻击者或内部恶意人员可能尝试修改、删除、插入或重排历史通信记录，以影响订单纠纷处理结果。如果系统仅保存普通数据库记录，而缺少完整性验证机制，则难以证明记录是否保持原始状态。

\item \textbf{管理员越权访问风险}：管理员具有较高审计权限，如果缺少权限控制和审计记录，管理员可能在无投诉、无异常或无授权的情况下查看敏感数据或执行溯源操作，造成二次隐私泄露。

\item \textbf{订单备注与标签泄露风险}：订单备注或标签可能包含用户生活习惯、地址细节、偏好信息等隐私内容。如果管理员端默认展示这些内容，可能导致审计功能扩大为隐私查看功能。
\end{enumerate}

针对上述威胁，FoodShield 采用多层防护机制。对于订单身份伪造问题，系统使用 HMAC-SM3 生成订单 Token，将订单号、PID、时间戳和随机数绑定到同一认证凭证中；对于 Token 重放风险，系统引入时间戳和随机数，并设置 Token 有效期；对于机器注册攻击，系统在注册阶段引入口令认证和设备识别，必要时结合注册频率限制；对于身份暴露风险，系统通过 PID 匿名身份机制隐藏真实用户身份，并限制骑手端和管理员端的可见字段；对于通信日志篡改问题，系统使用 SM3 消息摘要和 Merkle Tree 完整性验证机制检测历史记录是否被修改；对于管理员越权问题，系统要求管理员登录认证，并将关键操作写入审计日志；对于备注标签泄露问题，管理员端默认不展示用户备注或标签原文。

需要说明的是，本系统不追求对平台服务器的绝对匿名。平台服务器作为业务和安全控制中心，需要维护用户真实身份与 PID 的映射关系，以便在投诉、纠纷或审计场景下进行条件溯源。因此，FoodShield 采用的是“平台可控匿名”模型，而不是完全匿名模型。该模型的目标是在日常通信中对骑手、普通用户和外部攻击者隐藏真实身份，同时在满足审计条件时保留必要的责任追踪能力。
